\vspace*{\fill}
\begin{glyph}{Water (水)}{water}\label{glyph:water}
Water.
\end{glyph}

\begin{comps}
\comprtwocone & 水 \\\hline
\comprthreecone & Water (\ref{glyph:water}) \\\hline
\end{comps}

\begin{remglyph}
---\\\hline
\end{remglyph}

\begin{prons}
\pronsrtwocone & \textbf{スイ (SUI)} & \textbf{みず (mizu)}\\\hline
\pronsrthreecone & --- & --- \\\hline
\end{prons}

\begin{remprons}
\comkun{Make me zoom} on to the \hooglig{watery} \comon{chop suey}.\\\hline
\end{remprons}

\begin{somme}
水 & \textit{water} & みず (mizu)\\\hline
水中 & water + middle = \textit{underwater; in the water} & スイ{\middel}チュウ (SUI{\middel}CH\={U})\\\hline
下水 & lower + water = \textit{sewerage} & ゲ{\middel}スイ (GE{\middel}SUI)\\\hline
水玉 & water + jewel = \textit{water droplet} & みず{\middel}だま (mizu{\middel}dama)\\\hline
    \notyet{水星} & water + star = \textit{Mercury} & スイ{\middel}セイ (SUI{\middel}SEI)\\\hline
    \notyet{水曜日} & water + day of the week + sun (day) = \textit{Mercury} & スイ{\middel}ヨウ{\middel}び (SUI{\middel}Y\={O}{\middel}bi)\\\hline
    \notyet{水死} & water + death = \textit{drowning} & スイ{\middel}シ (SUI{\middel}SHI)\\\hline
\end{somme}

\begin{sentence}
\underline{あの} \underline{ダム}に\ruby{は}{}\ruby{\underline{水}}{みず}が\underline{ありますえん}。\\\hline
ano damu ni wa mizu ga arimasen.\\\hline
(that) (dam) (water) (isn't)\\\hline
=\textit{There is no water at that dam.}\\\hline
\end{sentence}
\end{multicols}
\vspace*{-\baselineskip}
\vhrulefill{5pt}
%%--------------------------------------------------------------------
%%--------------------------------------------------------------------
\begin{multicols}{2}
\begin{glyph}{Ease (安)}{ease}\label{glyph:ease}
Ease; stability; peacefulness.\\\hline
An important secondary meaning is ``inexpensive''.  In this regard, this kanji is the opposite of 高.
\end{glyph}

\begin{comps}
\comprtwocone & 宀 + 女 \\\hline
\comprthreecone & roof/cover + Woman (\ref{glyph:woman}) \\\hline
\end{comps}

\begin{remglyph}
Every \remcom{woman} needs a \remcom{roof} over her head for a sense of \hooglig{ease and stability}.\\\hline
\end{remglyph}

\begin{prons}
\pronsrtwocone & \textbf{アン (AN)} & \textbf{やす (yasu)}\\\hline
\pronsrthreecone & --- & --- \\\hline
\end{prons}

\begin{remprons}
\hooglig{Ease} your \comon{uncle} with \comkun{ya suit}.\\\hline
\end{remprons}

\begin{somme}
安い & \textit{easy; cheap} & やす{\middel}い (yasu{\middel}i)\\\hline
円安 & circle (yen) + ease = \textit{weak yen} & エン{\middel}だか (EN{\middel}daka)\\\hline
安心 & ease + heart = \textit{peace of mind} & アン{\middel}シン (AN{\middel}SHIN)\\\hline
安全 & ease + complete = \textit{safety} & アン{\middel}ゼン (AN{\middel}ZEN)\\\hline
    \notyet{不安} & not + ease = \textit{uneasiness} & フ{\middel}アン (FU{\middel}AN)\\\hline
    \notyet{公安} & public + ease = \textit{public peace} & コウ{\middel}アン (K\={O}{\middel}AN)\\\hline
\end{somme}

\begin{sentence}
    \ruby{\underline{安心}}{アンシン}{\;}\underline{して} \ruby{\underline{下}}{くだ}\underline{さい}。\\\hline
AN{\middel}SHIN shite kuda{\middel}sai.\\\hline
(peace of mind) (do) (please)\\\hline
=\textit{Please don't worry.}\\\hline
\end{sentence}
\end{multicols}
\vspace*{-\baselineskip}
\vhrulefill{5pt}
\vspace*{\fill}
%%--------------------------------------------------------------------
%%--------------------------------------------------------------------
\pagebreak
\vspace*{\fill}
\begin{multicols}{2}
\begin{glyph}{Strength (力)}{strength}\label{glyph:strength}
Strength; Power; Force.
\end{glyph}

\begin{comps}
\comprtwocone & 力 \\\hline
\comprthreecone & Strength (\ref{glyph:strength}) \\\hline
\end{comps}

\begin{remglyph}
---\\\hline
\end{remglyph}

\begin{prons}
\pronsrtwocone & \textbf{リョク、リキ (RYOKU, RIKI)} & \textbf{ちから (chikara)}\\\hline
\pronsrthreecone & --- & --- \\\hline
\rye{3}{リョク will be encountered far more often than リキ.  The kun-yomi is a common word on its own.}
\end{prons}

\begin{remprons}
By \hooglig{strengthening} the \comon{reeky}, \comon{lean yorkie}, its \comkun{cheek arrangement} will be re-established.\\\hline
\end{remprons}

\begin{somme}
力 & \textit{strenght} & ちから (chikara)\\\hline
全力 & complete + strength = \textit{all one's power} & ゼン{\middel}リョク (ZEN{\middel}RYOKU)\\\hline
有力 & have + strength = \textit{powerful} & ユウ{\middel}リョク (Y\={U}{\middel}RYOKU)\\\hline
水力 & water + strength = \textit{hydro power} & スイ{\middel}リョク (SUI{\middel}RYOKU)\\\hline
火力 & fire + strength = \textit{thermal power} & スイ{\middel}リョク (KA{\middel}RYOKU)\\\hline
    \notyet{電力} & electric + strength = \textit{electric power} & デン{\middel}リョク (DEN{\middel}RYOKU)\\\hline
    \notyet{力士} & strenght + gentleman = \textit{sumo wrestler} & リキ{\middel}シ (RIKI{\middel}SHI)\\\hline
\end{somme}

\begin{sentence}
    \ruby{\underline{中山}}{なかやま}\underline{さん}\ruby{は}{わ}\ruby{\underline{車}}{くるま}の\underline{ビジネス}で\underline{とても}{\;}\ruby{\underline{有力}}{ユウリョク}な\ruby{\underline{人}}{ひと}{\;}\underline{です}。\\\hline
Naka{\middel}yama-san wa kuruma no bijinesu de totemo Y\={U}{\middel}RYOKU na hito desu.\\\hline
(Nakayama-san) (car) (business) (very) (powerful) (person) (is)\\\hline
=\textit{In the automotive business, Nakayama-san is a very powerful person.}\\\hline
\end{sentence}
\end{multicols}
\vspace*{-\baselineskip}
\vhrulefill{5pt}
%%--------------------------------------------------------------------
%%--------------------------------------------------------------------
\begin{multicols}{2}
\begin{glyph}{Summer (夏)}{summer}\label{glyph:summer}
Summer.
\end{glyph}

\begin{comps}
\comprtwocone & 一 + 自 + 夂 \\\hline
\comprthreecone & One (\ref{glyph:one}) + Self (\ref{glyph:self}) + overtaking \\\hline
\end{comps}

\begin{remglyph}
    \hooglig{Summer} \remcom{overtakes} \remcom{one} \remcom{self}.\\\hline
\end{remglyph}

\begin{prons}
\pronsrtwocone & \textbf{カ (KA)} & \textbf{なつ (natsu)}\\\hline
\pronsrthreecone & ゲ (GE) & --- \\\hline
\end{prons}

\begin{remprons}
In \hooglig{summer} I \ucomon{get} \comkun{nuts} such that I start \comon{cussing}.\\\hline
\end{remprons}

\begin{somme}
夏 & \textit{summer} & なつ (natsu)\\\hline
    \notyet{夏休み} & summer + rest = \textit{summer vacation} & なつ　やす{\middel}み (natsu yasu{\middel}mi)\\\hline
    \notyet{夏物} & summer + thing = \textit{summer cloting} & なつ{\middel}もの (natsu{\middel}mono)\\\hline
    \notyet{立夏} & stand + summer = \textit{first day of summer} & リッ{\middel}カ (RIK{\middel}KA)\\\hline
\end{somme}

\begin{sentence}
\underline{この} \ruby{\underline{夏}}{なつ}に\ruby{は}{わ}\ruby{\underline{東京}}{トウキョウ}\ruby{へ}{え}\ruby{\underline{行}}{い}\underline{きました}。\\\hline
kono natsu ni wa T\={O}{\middel}KY\={O} e i{\middel}kimashite.\\\hline
(this) (summer) (Tokyo) (go)\\\hline
=\textit{(We'll) be going to Tokyo this summer.}\\\hline
\end{sentence}
\end{multicols}
\vspace*{-\baselineskip}
\vhrulefill{5pt}
\vspace*{\fill}
%%--------------------------------------------------------------------
%%--------------------------------------------------------------------
\pagebreak
\vspace*{\fill}
\begin{multicols}{2}
\begin{glyph}{Thousand (千)}{thousand}\label{glyph:thousand}
Thousand.\\\hline
Like 百, this kanji sometimes conveys a general idea of `many'.
\end{glyph}

\begin{comps}
\comprtwocone & ノ + 十 \\\hline
\comprthreecone & NO katakana/bend + Ten (\ref{glyph:thousand}) \\\hline
\end{comps}

\begin{remglyph}
    \remcom{No} \remcom{ten} \hooglig{thousand}.\\\hline
\end{remglyph}

\begin{prons}
\pronsrtwocone & \textbf{セン (SEN)} & ---\\\hline
\pronsrthreecone &  & いち (ichi) \\\hline
\end{prons}

\begin{remprons}
The \ucomkun{chief} is \comon{central} to the morale of \hooglig{thousands} of soldiers.\\\hline
\end{remprons}

\begin{somme}
千 & \textit{one thousand} & セン (SEN)\\\hline
二千 & two + thousand = \textit{two thousand} & ニ{\middel}セン (NI{\middel}SEN)\\\hline
三千 & three + thousand = \textit{three thousand} & サン{\middel}ゼン (SAN{\middel}ZEN)\\\hline
四千 & four + thousand = \textit{four thousand} & よん{\middel}セン (yon{\middel}SEN)\\\hline
五千 & five + thousand = \textit{five thousand} & ゴ{\middel}セン (GO{\middel}SEN)\\\hline
六千 & six + thousand = \textit{six thousand} & ロク{\middel}セン (ROKU{\middel}SEN)\\\hline
七千 & seven + thousand = \textit{seven thousand} & なな{\middel}セン (nana{\middel}SEN)\\\hline
八千 & eight + thousand = \textit{eight thousand} & ハチ{\middel}セン (HACHI{\middel}SEN)\\\hline
九千 & nine + thousand = \textit{nine thousand} & キュウ{\middel}セン (KY\={U}{\middel}SEN)\\\hline
\end{somme}

\begin{sentence}
    \ruby{\underline{米}}{こめ}\ruby{は}{わ}\ruby{\underline{七千円}}{ななセンエン}{\;}\underline{なので} \ruby{\underline{安}}{やす}\underline{くない}。\\\hline
kome wa nana{\middel}SEN{\middel}EN na no de yasu{\middel}kunai.\\\hline
(rice) (seven thousand yen) (because) (not cheap)\\\hline
=\textit{At seven thousand yen, rice isn't cheap.}\\\hline
\end{sentence}
\end{multicols}
\vspace*{-\baselineskip}
\vhrulefill{5pt}
%%--------------------------------------------------------------------
%%--------------------------------------------------------------------
\begin{multicols}{2}
\begin{glyph}{Rest (休)}{rest}\label{glyph:rest}
Rest.
\end{glyph}

\begin{comps}
\comprtwocone & 亻 + 木 \\\hline
\comprthreecone & Person (variant) (\ref{glyph:person}) + Tree (\ref{glyph:tree}) \\\hline
\end{comps}

\begin{remglyph}
\remcom{People} typically find \hooglig{rest} under \remcom{trees.}\\\hline
\end{remglyph}

\begin{prons}
\pronsrtwocone & \textbf{キュウ (KY\={U})} & \textbf{やす (yasu)}\\\hline
\pronsrthreecone & --- & --- \\\hline
\end{prons}

\begin{remprons}
If \comkun{ya suit} yourself to too much \hooglig{rest}, it will \comon{kill you}.\\\hline
\end{remprons}

\begin{somme}
休む & \textit{to rest} & やす{\middel}む (yasu{\middel}mu)\\\hline
お休み & \textit{``Good Night!''} & お{\middel}やす{\middel}み (o{\middel}yasu{\middel}mi)\\\hline
休日 & rest + day = \textit{holiday} & キュウ{\middel}ジツ (KY\={U}{\middel}JITSU)\\\hline
夏休み & summer + rest = \textit{summer vacation} & なつ　やす{\middel}み (natsu yasu{\middel}mi)\\\hline
冬休み & winter + rest = \textit{winter vacation} & ふゆ　やす{\middel}み (fuyu yasu{\middel}mi)\\\hline
休火山 & rest + fire + mountain = \textit{dormant volcano} & キュウ{\middel}カ{\middel}ザン (KY\={U}{\middel}KA{\middel}ZAN)\\\hline
\end{somme}

\begin{sentence}
    \ruby{\underline{山口}}{やまぐち}\underline{さん}\ruby{は}{わ}\ruby{\underline{半日}}{ハンニチ}{\;}\ruby{\underline{休}}{やす}\underline{みました}。\\\hline
Yama{\middel}guchi-san wa HAN{\middel}NICHI yasu{\middel}mimashita.\\\hline
(Yamaguchi-san) (half a day) (rested)\\\hline
=\textit{Yamaguchi-san rested for half a day.}\\\hline
\end{sentence}
\end{multicols}
\vspace*{-\baselineskip}
\vhrulefill{5pt}
\vspace*{\fill}
%%--------------------------------------------------------------------
%%--------------------------------------------------------------------
\pagebreak
\vspace*{\fill}
\begin{multicols}{2}
\begin{glyph}{Hand (手)}{te}\label{glyph:hand}
Hand.
\end{glyph}

\begin{comps}
\comprtwocone & 手 \\\hline
\comprthreecone & Hand (\ref{glyph:hand}) \\\hline
\end{comps}

\begin{remglyph}
---\\\hline
\end{remglyph}

\begin{prons}
\pronsrtwocone & \textbf{シュ (SHU)} & \textbf{て (te)}\\\hline
\pronsrthreecone & --- & た (ta) \\\hline
\end{prons}

\begin{remprons}
\hooglig{Hand} me the \comkun{tempting} \ucomkun{tunnel} \comon{shooter}.\\\hline
\end{remprons}

\begin{somme}
手 & \textit{hand} & て (te)\\\hline
右手 & right + hand = \textit{right hand} & みぎ{\middel}て (migi{\middel}te)\\\hline
左手 & left + hand = \textit{left hand} & ひだり{\middel}て (hidari{\middel}te)\\\hline
手元 & hand + basis = \textit{at hand} & て{\middel}もと (te{\middel}moto)\\\hline
手首 & hand + neck = \textit{wrist} & て{\middel}くび (te{\middel}kubi)\\\hline
    \notyet{手話} & hand + speak = \textit{sign language} & シュ{\middel}ワ (SHU{\middel}WA)\\\hline
    \notyet{空手} & empty + hand = \textit{karate} & から{\middel}て (kara{\middel}te)\\\hline
\end{somme}

\begin{irrr}
上手 & upper + hand = \textit{to be good at (something)} & ジョウ{\middel}ず (J\={O}{\middel}zu)\\\hline
下手 & lower + hand = \textit{to be poor at (something)} & へ{\middel}た (he{\middel}ta)\\\hline
\end{irrr}

\begin{sentence}
\underline{あの} \underline{ピアニスト}の\ruby{\underline{手}}{て}\ruby{は}{わ}\ruby{\underline{美}}{うつく}\underline{しい} \underline{ですね}。\\\hline
ano pianisuto no te wa utsuku{\middel}shii desu ne.\\\hline
(that) (pianist) (hands) (beautiful) (aren't they)\\\hline
=\textit{That pianist's hands are beautiful, aren't they.}\\\hline
\end{sentence}
\end{multicols}
\vspace*{-\baselineskip}
\vhrulefill{5pt}
%%--------------------------------------------------------------------
%%--------------------------------------------------------------------
\begin{multicols}{2}
\begin{glyph}{Main (本)}{main}\label{glyph:main}
Main; Origin.  Also used to emphasize the idea of ``this''.\\\hline
When encountered alone, it often indicates a book.\\\hline
It is also a `counter' for long, cylindrical objects such as pencils and bottles.\\\hline
The sense of ``origin'', incidentally, is evident in the compound for Japan itself, the ``Land of the Rising Sun.''
\end{glyph}

\begin{comps}
\comprtwocone & 木 + 一 \\\hline
\comprthreecone & Tree (\ref{glyph:tree}) + One (\ref{glyph:one}) \\\hline
\end{comps}

\begin{remglyph}
One main tree.\\\hline
\end{remglyph}

\begin{prons}
\pronsrtwocone & \textbf{ホン (HON)} & \textbf{もと (moto)}\\\hline
\pronsrthreecone & --- & --- \\\hline
\rye{3}{もと is common in Japanese family and place names.}
\end{prons}

\begin{remprons}
The \comkun{motorists} were \hooglig{mainly} \comon{haunted}.\\\hline
\end{remprons}

\begin{somme}
本 & \textit{book} & ホン (HON)\\\hline
日本 & sun + main = \textit{Japan} & ニ{\middel}ホン、ニッ{\middel}ポン (NI{\middel}HON, NIP{\middel}PON)\\\hline
全日本 & complete + sun + main = \textit{all Japan} & ゼン{\middel}ニ{\middel}ホン (ZEN{\middel}NI{\middel}HON)\\\hline
日本人 & sun + main + person = \textit{a Japanese person} & ニ{\middel}ホン{\middel}ジン (NI{\middel}HON{\middel}JIN)\\\hline
山本さん & mountain + main = \textit{Yamamoto-san} & やま{\middel}もと{\middel}さん (yama{\middel}moto{\middel}san)\\\hline
    \notyet{日本語} & sun + main + words = \textit{Japanese (language)} & ニ{\middel}ホン{\middel}ゴ (NI{\middel}HON{\middel}GO)\\\hline
    \notyet{本土} & main + earth = \textit{mainland} & ホン{\middel}ド (HON{\middel}DO)\\\hline
\end{somme}

\begin{sentence}
\ruby{\underline{山下}}{やました}\underline{さん}\ruby{は}{わ}\ruby{\underline{高}}{たか}\underline{い} \ruby{\underline{本}}{ホン}\ruby{を}{お}\ruby{\underline{買}}{か}\underline{いました}。\\\hline
Yama{\middel}shita-san wa taka{\middel}i HON o ka{\middel}imashita.\\\hline
(Yamashita-san) (expensive) (book) (bought)\\\hline
=\textit{Yamashita-san bought an expensive book.}\\\hline
\end{sentence}
\end{multicols}
\vspace*{-\baselineskip}
\vhrulefill{5pt}
\vspace*{\fill}
%%--------------------------------------------------------------------
%%--------------------------------------------------------------------
\pagebreak
\vspace*{\fill}
\begin{multicols}{2}
\begin{glyph}{Change (化)}{change}\label{glyph:change}
Change; Transform.\\\hline
When used as a suffix, this character expresses a variety of English words ending in ``-ification'' or ``-ization''.
\end{glyph}

\begin{comps}
\comprtwocone & 亻 + 匕 \\\hline
\comprthreecone & Person (variant) (\ref{glyph:person}) + spoon \\\hline
\end{comps}

\begin{remglyph}
That \remcom{person} \hooglig{changed} the \remcom{spoon}.\\\hline
\end{remglyph}

\begin{prons}
\pronsrtwocone & \textbf{カ (KA)} & ---\\\hline
\pronsrthreecone & KE (ケ) & ば (ba) \\\hline
\end{prons}

\begin{remprons}
People \comon{cussed} the \ucomkun{bulging} \ucomon{keg} that \hooglig{changed} the alcoholic taste.\\\hline
\end{remprons}

\begin{somme}
国有化 & country + have + change = \textit{nationalization} & コク{\middel}ユウ{\middel}カ (KOKU{\middel}Y\={U}{\middel}KA)\\\hline
    \notyet{美化} & beautiful + change = \textit{beautification} & ビ{\middel}カ (BI{\middel}KA)\\\hline
    \notyet{化学} & change + study = \textit{chemistry} & キ{\middel}ガク (KA{\middel}GAKU)\\\hline
    \notyet{気化} & spirit + change = \textit{vaporization} & キ{\middel}カ (KI{\middel}KA)\\\hline
    \notyet{電化} & electric + change = \textit{electrification} & デン{\middel}カ (DEN{\middel}KA)\\\hline
    \notyet{同化} & same + change = \textit{assimilation} & ドウ{\middel}カ (D\={O}{\middel}KA)\\\hline
\end{somme}

\begin{sentence}
\ruby{\underline{山本}}{やまもと}\underline{さん}\ruby{は}{わ}\ruby{\underline{化学}}{カガク}が\ruby{\underline{大好}}{ダイす}\underline{き}です。\\\hline
Yama{\middel}moto-san wa KA{\middel}GAKU ga DAI su{\middel}ki desu.\\\hline
(Yamamoto-san) (chemistry) (really likes)\\\hline
=\textit{Yamamoto-san really likes chemistry.}\\\hline
\end{sentence}
\end{multicols}
\vspace*{-\baselineskip}
\vhrulefill{5pt}
%%--------------------------------------------------------------------
%%--------------------------------------------------------------------
\begin{multicols}{2}
\begin{glyph}{River (川)}{river}\label{glyph:river}
River.
\end{glyph}

\begin{comps}
\comprtwocone & 川 \\\hline
\comprthreecone & River (\ref{glyph:river}) \\\hline
\end{comps}

\begin{remglyph}
---\\\hline
\end{remglyph}

\begin{prons}
\pronsrtwocone & --- & \textbf{カワ (kawa)}\\\hline
\pronsrthreecone & セン (SEN) & --- \\\hline
\rye{3}{When appearing as the final character in the names of rivers, this kanji is invariable voiced as がわ.}
\end{prons}

\begin{remprons}
The \ucomon{central} \hooglig{river} attraction was the \comkun{carwash.}\\\hline
\end{remprons}

\begin{somme}
川 & \textit{river} & かわ (kawa)\\\hline
川上 &  river + upper = \textit{upstream} & かわ{\middel}かみ (kawa{\middel}kami)\\\hline
川下 & river + lower = \textit{downstream} & かわ{\middel}した (kawa{\middel}shita)\\\hline
川口 & river + mouth = \textit{river mouth} & かわ{\middel}ぐち (kawa{\middel}guchi)\\\hline
川本さん & river + main = \textit{Kawamoto-san} & かわ{\middel}もと{\middel}さん (kawa{\middel}moto{\middel}san)\\\hline
\end{somme}

\begin{sentence}
\ruby{\underline{山}}{やま}の\ruby{\underline{上}}{うえ}に\ruby{\underline{行}}{い}\underline{く}と\ruby{\underline{川}}{かわ}が\ruby{\underline{見}}{み}\underline{えます}か。\\\hline
yama no ue ni i{\middel}ku to kawa ga mi{\middel}emasu ka.\\\hline
(mountain) (upper) (go) (river) (can see)\\\hline
=\textit{Can you see the river if you go to the top of the mountain.}\\\hline
\end{sentence}
\end{multicols}
\vspace*{-\baselineskip}
\vhrulefill{5pt}
\vspace*{\fill}
%%--------------------------------------------------------------------
%%--------------------------------------------------------------------
\pagebreak
\vspace*{\fill}
\begin{multicols}{2}
\begin{glyph}{Body (体)}{body}\label{glyph:body}
The physical body of a person or object.\\\hline
A secondary meaning of ``substance'' and ``style''.
\end{glyph}

\begin{comps}
\comprtwocone & 亻 + 本 \\\hline
\comprthreecone & Person (variant) (\ref{glyph:person}) + Main (\ref{glyph:main}) \\\hline
\end{comps}

\begin{remglyph}
The \hooglig{body} is the \remcom{main} aspect of any \remcom{person}.\\\hline
\end{remglyph}

\begin{prons}
\pronsrtwocone & \textbf{タイ (TAI)} & \textbf{からだ (karada)}\\\hline
\pronsrthreecone & テイ (TEI) & --- \\\hline
\end{prons}

\begin{remprons}
My \hooglig{body} \comon{tightened} as the \comkun{car adapter} got \ucomon{taped} to me.\\\hline
\end{remprons}

\begin{somme}
体 & \textit{body} & からだ (karada)\\\hline
大体 & large + body = \textit{in general} & ダイ{\middel}タイ (DAI{\middel}TAI)\\\hline
全体 & complete + body = \textit{the whole} & ゼン{\middel}タイ (ZEN{\middel}TAI)\\\hline
肉体 & meat + body = \textit{the body} & ニク{\middel}タイ (NIKU{\middel}TAI)\\\hline
体力 & body + strength = \textit{physical strength} & タイ{\middel}リョク (TAI{\middel}RYOKU)\\\hline
自体 & self + body = \textit{in itself} & ジ{\middel}タイ (JI{\middel}TAI)\\\hline
    \notyet{団体} & group + body = \textit{organization} & ダン{\middel}タイ (DAN{\middel}TAI)\\\hline
\end{somme}

\begin{sentence}
\underline{あの} \ruby{\underline{人}}{ひと}\ruby{は}{わ}\underline{すごく} \ruby{\underline{体力}}{タイリョク}が\ruby{\underline{有}}{あ}\underline{ります}。\\\hline
ano hito wa sugoku TAI{\middel}RYOKU ga a{\middel}rimasu.\\\hline
(that) (person) (incredible) (physical strength) (has)\\\hline
=\textit{That person has incredible physical strength.}\\\hline
\end{sentence}
\end{multicols}
\vspace*{-\baselineskip}
\vhrulefill{5pt}
%%--------------------------------------------------------------------
%%--------------------------------------------------------------------
\begin{multicols}{2}
\begin{glyph}{North (北)}{north}\label{glyph:north}
North.
\end{glyph}

\begin{comps}
\comprtwocone & 爿 + 匕 \\\hline
\comprthreecone & bed + spoon \\\hline
\end{comps}

\begin{remglyph}
    The \hooglig{northern} \remcom{bed} \remcom{spoon}.\\\hline
\end{remglyph}

\begin{prons}
\pronsrtwocone & \textbf{ホク (HOKU)} & \textbf{きた (kita)}\\\hline
\pronsrthreecone & --- & --- \\\hline
\end{prons}

\begin{remprons}
\comon{Hawks} move \hooglig{northward} to the \comkun{citadel}.\\\hline
\end{remprons}

\begin{somme}
北 & \textit{north} & きた (kita)\\\hline
北米 & north + rice = \textit{North America} & ホク{\middel}ベイ (HOKU{\middel}BEI)\\\hline
北口 & north + mouth = \textit{north exit/entrance} & きた{\middel}ぐち (kita{\middel}guci)\\\hline
    \notyet{北東} & north + east = \textit{northeast} & ホク{\middel}トウ (HOKU{\middel}T\={O})\\\hline
    \notyet{北西} & north + west = \textit{northwest} & ホク{\middel}セイ (HOKU{\middel}SEI)\\\hline
    \notyet{北海道} & north + sea + road = \textit{Hokkaido} & ホッ{\middel}カイ{\middel}ドウ (HOK{\middel}KAI{\middel}D\={O})\\\hline
\end{somme}

\begin{sentence}
\ruby{\underline{北日本}}{きたニホン}の\ruby{\underline{冬}}{ふゆ}\ruby{は}{わ}\underline{とても} \ruby{\underline{美}}{うつく}\underline{しい}。\\\hline
kita{\middel}NI{\middel}HON no fuyu wa totemo utsuku{\middel}shii.\\\hline
(northern Japan) (winter) (very) (beautiful)\\\hline
=\textit{Northern Japan's winter is very beautiful.}\\\hline
\end{sentence}
\end{multicols}
\vspace*{-\baselineskip}
\vhrulefill{5pt}
\vspace*{\fill}
%%--------------------------------------------------------------------
%%--------------------------------------------------------------------
\pagebreak
\vspace*{\fill}
\begin{multicols}{2}
\begin{glyph}{Rice field (田)}{rice_field}\label{glyph:rice_field}
Rice field, or a field in general.
\end{glyph}

\begin{comps}
\comprtwocone & 田 \\\hline
\comprthreecone & Rice Field (\ref{glyph:rice_field}) \\\hline
\end{comps}

\begin{remglyph}
---\\\hline
\end{remglyph}

\begin{prons}
\pronsrtwocone & \textbf{デン (DEN)} & \textbf{た (ta)}\\\hline
\pronsrthreecone & --- & --- \\\hline
\rye{3}{This kanji is commonly used in Japanese family names; it is often voiced, as in the fourth and fifth examples below, the latter of which is a mix of on and kun-yomi.}
\end{prons}

\begin{remprons}
In the \hooglig{rice field} was a \comkun{tub} hid inside a \comon{den}.\\\hline
\end{remprons}

\begin{somme}
田 & \textit{rice field} & た (ta)\\\hline
水田 & water + rice field = \textit{rice paddy} & スイ{\middel}デン (SUI{\middel}DEN)\\\hline
田中さん & rice field + middle = \textit{Tanaka-san} & た{\middel}なか{\middel}さん (ta{\middel}naka{\middel}san)\\\hline
山田さん & mountain + rice field = \textit{Yamada-san} & やま{\middel}だ{\middel}さん (yama{\middel}da{\middel}san)\\\hline
本田さん & main + rice field = \textit{Honda-san} & ホン{\middel}だ{\middel}さん (HON{\middel}da{\middel}san)\\\hline
\end{somme}

\begin{irrr}
    田舎{\notcov} & rice field + building = \textit{the countryside} & いなか (inaka)\\\hline
\end{irrr}

\begin{sentence}
\ruby{\underline{水田}}{スイデン}\ruby{を}{お}\ruby{\underline{見}}{み}に\ruby{\underline{行}}{い}\underline{きました}。\\\hline
SUI{\middel}DEN o mi ni i{\middel}kimashita.\\\hline
(rice paddy) (see) (went)\\\hline
=\textit{(She) went to see the rice paddies.}\\\hline
\end{sentence}
\end{multicols}
\vspace*{-\baselineskip}
\vhrulefill{5pt}
%%--------------------------------------------------------------------
%%--------------------------------------------------------------------
\begin{multicols}{2}
\begin{glyph}{Man (男)}{man}\label{glyph:man}
Man; Male.
\end{glyph}

\begin{comps}
\comprtwocone & 田 + 力 \\\hline
\comprthreecone & Rice Field (\ref{glyph:rice_field}) + Strength (\ref{glyph:strength})\\\hline
\end{comps}

\begin{remglyph}
A real \hooglig{man} shows his \remcom{strength} in the \remcom{rice fields}.\\\hline
\end{remglyph}

\begin{prons}
\pronsrtwocone & \textbf{ダン (DAN)} & \textbf{おとこ (otoko)}\\\hline
\pronsrthreecone & ナン (NAN) & --- \\\hline
\end{prons}

\begin{remprons}
\hooglig{Man} perfected the \comkun{auto-complete} function to dunk \ucomon{nuns}.\\\hline
\end{remprons}

\begin{somme}
男 & \textit{man} & おとこ (otoko)\\\hline
男女 & man + woman = \textit{men and woman} & ダン{\middel}ジョ (DAN{\middel}JO)\\\hline
男子 & man + child = \textit{men's} & ダン{\middel}シ (DAN{\middel}SHI)\\\hline
大男 & large + man = \textit{large man} & おお{\middel}おとこ (\={o}{\middel}otoko)\\\hline
    \notyet{雪男} & snow + man = \textit{the abominable snowman} & ゆき{\middel}おとこ (yuki{\middel}otoko)\\\hline
\end{somme}

\begin{sentence}
\underline{あの} \ruby{\underline{男}}{おとこ}の\ruby{\underline{人}}{ひと}\ruby{は}{わ}\ruby{\underline{五千円}}{ゴセンエン}\ruby{を}{お}\underline{なくした}。\\\hline
ano otoko no hito wa GO{\middel}SEN{\middel}EN o nakushita.\\\hline
(that) (man) (person) (five thousand yen) (lost)\\\hline
=\textit{That man lost five thousand yen.}\\\hline
\end{sentence}
\end{multicols}
\vspace*{-\baselineskip}
\vhrulefill{5pt}
\vspace*{\fill}
%%--------------------------------------------------------------------
%%--------------------------------------------------------------------
\pagebreak
\vspace*{\fill}
\begin{multicols}{2}
\begin{glyph}{House (家)}{house}\label{glyph:house}
Some sense of relation to house or family.\\\hline
When used as a suffix, it denotes a person having some degree of skill or interest with respect to the kanji preceding it.
\end{glyph}

\begin{comps}
\comprtwocone & 宀 + 豕 \\\hline
\comprthreecone & roof + pig\\\hline
\end{comps}

\begin{remglyph}
    A \hooglig{house} is a \remcom{roof} over a \remcom{pig}.\\\hline
\end{remglyph}

\begin{prons}
\pronsrtwocone & \textbf{カ (KA)} & \textbf{いえ (ie)}\\\hline
\pronsrthreecone & ケ (KE) & や (ya) \\\hline
\end{prons}

\begin{remprons}
\comkun{i.e.} we started \comon{cussing} once we realised the \ucomon{keg} was already \hooglig{housed} in the \ucomkun{yacht}.\\\hline
\end{remprons}

\begin{somme}
家 & \textit{house} & いえ (ie)\\\hline
人家 & person + house = \textit{dwelling} & ジン{\middel}カ (JIN{\middel}KA)\\\hline
    \notyet{売り家} & sell + house = \textit{house for sale} & う{\middel}り　いえ (u{\middel}ri ie)\\\hline
    \notyet{家具} & house + tool = \textit{furniture} & カ{\middel}グ (KA{\middel}GU)\\\hline
    \notyet{愛鳥家} & love + bird + house = \textit{bird lover} & アイ{\middel}チョウ{\middel}カ (AI{\middel}CH\={O}{\middel}KA)\\\hline
\end{somme}

\begin{sentence}
\ruby{\underline{本田}}{ホンだ}\underline{さん}の\ruby{\underline{家}}{いえ}\ruby{を}{お}\ruby{\underline{見}}{み}\underline{ました}か。\\\hline
HON{\middel}da{\middel}san no ie o mi{\middel}mashita ka.\\\hline
(Honda-san) (house) (see)\\\hline
=\textit{Did you see Honda-san's house?}\\\hline
\end{sentence}
\end{multicols}
\vspace*{-\baselineskip}
\vhrulefill{5pt}
%%--------------------------------------------------------------------
%%--------------------------------------------------------------------
\begin{multicols}{2}
\begin{glyph}{East (東)}{east}\label{glyph:east}
East.
\end{glyph}

\begin{comps}
\comprtwocone & 木 + 日 \\\hline
\comprthreecone & Tree (\ref{glyph:tree}) + Sun (\ref{glyph:tree})\\\hline
\end{comps}

\begin{remglyph}
In the \hooglig{east} the \remcom{sun} rises behind the \remcom{trees}.\\\hline
\end{remglyph}

\begin{prons}
\pronsrtwocone & \textbf{トウ (T\={O})} & \textbf{ひがし (higashi)}\\\hline
\pronsrthreecone & --- & --- \\\hline
\end{prons}

\begin{remprons}
\hooglig{Eastern} \comon{torture} methods will make \comkun{him gush}.\\\hline
\end{remprons}

\begin{somme}
東 & \textit{east} & ひがし (higashi)\\\hline
東口 & east + mouth = \textit{east exit/entrance} & ひがし{\middel}ぐち (higashi{\middel}guchi)\\\hline
中東 & middle + east = \textit{The Middle East} & チュウ{\middel}トウ (NAN{\middel}T\={O})\\\hline
東京 & east + capital = \textit{Tokyo} & キョウ{\middel}トウ (KY\={O}{\middel}T\={O})\\\hline
南東 & south + east = \textit{south east} & ナン{\middel}トウ (NAN{\middel}T\={O})\\\hline
北北東 & north + north + east = \textit{north-northeast} & ホク{\middel}ホク{\middel}トウ (HOKU{\middel}HOKU{\middel}T\={O})\\\hline
\end{somme}

\begin{sentence}
\ruby{\underline{東口}}{ひがしぐち}の\underline{キオスク}で\ruby{\underline{本}}{ホン}\ruby{を}{お}\ruby{\underline{買}}{か}\underline{いました}。\\\hline
higashi{\middel}guchi no kiosuku de HON o ka{\middel}imashita.\\\hline
(east exit) (kiosk) (book) (bought)\\\hline
=\textit{(I) bought a book at the east exit kiosk.}\\\hline
\end{sentence}
\end{multicols}
\vspace*{-\baselineskip}
\vhrulefill{5pt}
\vspace*{\fill}
%--------------------------------------------------------------------
%--------------------------------------------------------------------
\pagebreak
\vspace*{\fill}
\begin{multicols}{2}
\begin{glyph}{Think (思)}{think}\label{glyph:think}
Think.
\end{glyph}

\begin{comps}
\comprtwocone & 田 + 心 \\\hline
\comprthreecone & Rice field (\ref{glyph:rice_field}) + Heart (\ref{glyph:heart})\\\hline
\end{comps}

\begin{remglyph}
I \hooglig{think} the \remcom{rice field} has forever stolen my \remcom{heart}.\\\hline
\end{remglyph}

\begin{prons}
\pronsrtwocone & \textbf{シ (SHI)} & \textbf{おも (omo)} \\\hline
\pronsrthreecone & --- & --- \\\hline
\end{prons}

\begin{remprons}
I \hooglig{think} \comkun{omo} will give it a nice \comon{sheen}.\\\hline
\end{remprons}

\begin{somme}
思う & \textit{to think} & おも{\middel}う (omo{\middel}u)\\\hline
    \notyet{思い出す} & think + exit = \textit{to remember} & おも{\middel}い　だ{\middel}す (omo{\middel}i da{\middel}su)\\\hline
    \notyet{物思い} & thing + think = \textit{pensiveness} & もの　おも{\middel}い (mono omo{\middel}i)\\\hline
    \notyet{思い切る} & think + cut = \textit{one's intent} & おも{\middel}い　き{\middel}る (omo{\middel}i ki{\middel}ru)\\\hline
    \notyet{意思} & mind + think = \textit{(one's) intent} & イ{\middel}シ (I{\middel}SHI)\\\hline
\end{somme}

\begin{sentence}
    \underline{テーベル}の\ruby{\underline{上}}{うえ}に\ruby{\underline{四百円}}{よんヒャク}{\;}\underline{ある}と\ruby{\underline{思}}{おも}\underline{います}。\\\hline
t\={e}beru no ue ni yon{\middel}HYAKU{\middel}EN aru to omo{\middel}imasu.\\\hline
(table) (above) (four hundred yen) (is) (think)\\\hline
=\textit{(I) think there's four hundred yen on the table.}\\\hline
\end{sentence}
\end{multicols}
\vspace*{-\baselineskip}
\vhrulefill{5pt}
%%--------------------------------------------------------------------
%%--------------------------------------------------------------------
\begin{multicols}{2}
\begin{glyph}{Ear (耳)}{ear}\label{glyph:ear}
Ear.
\end{glyph}

\begin{comps}
\comprtwocone & 耳 \\\hline
\comprthreecone & Ear (\ref{glyph:ear})\\\hline
\end{comps}

\begin{remglyph}
---\\\hline
\end{remglyph}

\begin{prons}
\pronsrtwocone & --- & \textbf{みみ (mimi)} \\\hline
\pronsrthreecone & ジ (JI) & --- \\\hline
\end{prons}

\begin{remprons}
\comkun{Mimic} the \hooglig{ear} \ucomon{jig}.\\\hline
\end{remprons}

\begin{somme}
耳 & \textit{ear} & 耳 (mimi)\\\hline
耳元 & ear + basis = \textit{close to one's ear} & みみ{\middel}もと (mimi{\middel}moto)\\\hline
    \notyet{空耳} & empty + ear = \textit{mishearing; feigned deafness} & そら{\middel}みみ (sora{\middel}mimi)\\\hline
\end{somme}

\begin{sentence}
\ruby{\underline{田中}}{たなか}\underline{さん}の\ruby{\underline{耳}}{みみ}が\ruby{\underline{赤}}{あか}\underline{い}。\\\hline
Ta{\middel}naka-san no mimi ga aka{\middel}i.\\\hline
(Tanaka-san) (ears) (red)\\\hline
=\textit{Tanaka-san's ears are red.}\\\hline
\end{sentence}
\end{multicols}
\vspace*{-\baselineskip}
\vhrulefill{5pt}
\vspace*{\fill}
%%--------------------------------------------------------------------
%%--------------------------------------------------------------------
\pagebreak
\vspace*{\fill}
\begin{multicols}{2}
\begin{glyph}{Father (父)}{father}\label{glyph:father}
Father.
\end{glyph}

\begin{comps}
\comprtwocone & 父 \\\hline
\comprthreecone & Father (\ref{glyph:father})\\\hline
\end{comps}

\begin{remglyph}
---\\\hline
\end{remglyph}

\begin{prons}
\pronsrtwocone & \textbf{フ (FU)} & \textbf{ちち (chichi)} \\\hline
\pronsrthreecone & --- & --- \\\hline
\end{prons}

\begin{remprons}
Any \comkun{chichi} \hooglig{father} is a \comon{foofy}.\\\hline
\end{remprons}

\begin{somme}
\rye{3}{The various words for father shown below reflect differing levels of familiarity and politeness; most terms used for addressing people in Japan have multiple forms such as these.}
父 & \textit{father} & ちち (chichi)\\\hline
父の日 & father + sun (day) = \textit{Father's Day} & ちち{\middel}の{\middel}ひ (chichi{\middel}no{\middel}oya)\\\hline
    \notyet{父母} & father + mother = \textit{father and mother} & フ{\middel}ボ (FU{\middel}BO)\\\hline
    \notyet{父親} & fater + parent = \textit{father} & ちち{\middel}おや (chichi{\middel}oya)\\\hline
\end{somme}

\begin{irrr}
お父さん & \textit{father} & おとうさん (ot\={o}san)\\\hline
    伯父さん{\notcov} & related + father = \textit{uncle} & おじさん (ojisan)\\\hline
\end{irrr}

\begin{sentence}
\ruby{\underline{父}}{ちち}\ruby{は}{わ}\underline{ビジネス}で\ruby{\underline{東京}}{トウキョウ}\ruby{へ}{え}\ruby{\underline{行}}{い}\underline{きました}。\\\hline
chichi wa bijinesu de T\={O}{\middel}KY\={O} e i{\middel}kimashita.\\\hline
(Father) (business) (Tokyo) (went)\\\hline
=\textit{Father went to Tokyo on business.}\\\hline
\end{sentence}
\end{multicols}
\vspace*{-\baselineskip}
\vhrulefill{5pt}
%%--------------------------------------------------------------------
%%--------------------------------------------------------------------
\begin{multicols}{2}
\begin{glyph}{Say (言)}{say}\label{glyph:say}
``Say'' and things of a verbal nature.\\\hline
This character appears in more than 60 kanji.
\end{glyph}

\begin{comps}
\comprtwocone & 言 \\\hline
\comprthreecone & Speech/Speaking/Say (\ref{glyph:say})\\\hline
\end{comps}

\begin{remglyph}
Sound waves traveling outward.\\\hline
\end{remglyph}

\begin{prons}
\pronsrtwocone & \textbf{ゲン (GEN)} & \textbf{い (i)} \\\hline
\pronsrthreecone & ゴン (GON) & こと (koto) \\\hline
\end{prons}

\begin{remprons}
They \hooglig{say} \comon{Genghis Khan's} \comkun{interest} in \ucomkun{koto's} is \ucomon{gone} with the wind.\\\hline
\end{remprons}

\begin{somme}
言う & \textit{to say} & い{\middel}う (i{\middel}u)\\\hline
言い回し & say + rotate = \textit{turn of expression} & い{\middel}い　まわ{\middel}し (i{\middel}i mawa{\middel}shi)\\\hline
明言 & bright + say = \textit{declaration} & メイ{\middel}ゲン (MEI{\middel}GEN)\\\hline
    \notyet{言語} & say + words = \textit{language} & ゲン{\middel}ゴ (GEN{\middel}GO)\\\hline
    \notyet{言い出す} & say + exit = \textit{to begin to say} & い{\middel}い　だ{\middel}す (i{\middel}i da{\middel}su)\\\hline
    \notyet{言い切る} & say + cut = \textit{to say definitively} & い{\middel}い　き{\middel}る (i{\middel}i ki{\middel}ru)\\\hline
    言葉{\notcov} & say + leaf = \textit{word} & こと{\middel}ば (koto{\middel}ba)\\\hline
\end{somme}

\begin{sentence}
\ruby{\underline{川本}}{かわもと}\underline{さん} \underline{も} \underline{``はい''}と\ruby{\underline{言}}{い}\underline{いました}。\\\hline
Kawa{\middel}moto-san mo ``hai'' to i{\middel}imashita.\\\hline
(Kawamoto-san) (also) (``yes'') (said)\\\hline
=\textit{Kawamoto-san also said ``yes''.}\\\hline
\end{sentence}
\end{multicols}
\vspace*{-\baselineskip}
\vhrulefill{5pt}
\vspace*{\fill}
%%--------------------------------------------------------------------
%%--------------------------------------------------------------------
\pagebreak
